% Volume VI: Adversaries, Platforms, and Automated Judgment
% Part XI. Untrustworthy Actors
% Chapter 62: Informants, Whistleblowers, and Anonymous Speech
% Spine: proof-spines/ch062-spine.md

\chapter{Informants, Whistleblowers, and Anonymous Speech}
\label{ch:ch062}


\textbf{Anonymity} can protect corrective information while also reducing accountability. The outline therefore refuses both easy romanticism and easy suspicion. Anonymous speech matters because retaliation, surveillance, and captured channels often make attribution incompatible with disclosure; yet the loss of identity information shifts the burden of trust onto other features of the record.

The chapter treats anonymity as a redistribution of \emph{verification burdens}. A nameless allegation is neither self-authenticating nor self-discrediting. What matters is whether the claim can be corroborated, whether documents or observations can be independently checked, and whether the surrounding pattern of evidence survives scrutiny without leaning on identity alone. Anonymous speech is thus a protective condition for some truths to surface, but never a substitute for verification.
